\chapter{Conclusion}
\label{ch:conclusion}

The mathematical legitimacy of an inhibitory point process, identification of its observation law and predictive value of a fitted recipe are separate achievements. The present contribution connects them through restricted stability and identification results, controlled nuisance-sensitivity comparisons and exploratory prediction in a reconstructed earthquake catalogue.

The finite-mode physical model supports first-moment stability under explicit sign, mark, state-domain and coupling assumptions. Infinite stationary second moments can coexist with that result. Finite-horizon approximation does not exchange infinite-mode and stationary limits, and the no-feedback exact-memory counterexamples establish an obstruction without settling active nonlinear inhibition. At zero feedback the diagnostic stress kernel disappears from the observation law. Positive-feedback identification is possible under the stated full-law assumptions, including known transforms and rich mark support; it is not a guarantee of stable finite-catalogue estimation. Conditional amplitude profiling and fixed-rule score response provide further finite-history tools within those boundaries.

The controlled experiments change the interpretation of weak fitted detection. The known-law procedure detected 150/150 saved alternatives, with 50/1,000 separate null rejections. Background level estimation retained 123/150 detections when trend was supplied; estimating level and trend retained 68/150. These are pronounced procedure differences under the declared law. The oracle receives additional nuisance information, so the contrast does not separate unavoidable composite-testing difficulty from recoverable procedure inefficiency. Fixing unknown real-data backgrounds merely to reproduce the synthetic advantage would be unjustified.

In the observed comparison, the shared-inhibition recipe improved score in none of three declared periods, while an Omori triggering comparator scored better than Mittag--Leffler triggering. Approximately 93.7\% of the inhibition recipe's aggregate disadvantage occurred in the first period; one later contrast was almost zero. Reused periods, overlapping training prefixes, unequal kernel flexibility and residual adequacy concerns restrict the result to exploratory recipe comparisons. It is useful negative predictive evidence, without establishing the presence or absence of physical inhibition or shared material memory.

The aggregation calculation gives a substantive reason to investigate locality: equal pooled depletion can conceal different aggregate local rates. A bounded comparison of no inhibition, pooled inhibition and local inhibition would need common observations, matched backgrounds, serious triggering benchmarks and complete regional compensators. This remains an extension rather than a condition for the existing contribution. The next empirical question is whether locality adds predictive information after background and triggering structure have been matched, rather than whether another search can produce a preferred sign.
